\section{Экспериментальный план и постановка замеров}

\DTLloadrawdb{nlpparams}{tables/nlp_illustrative_params.csv}%
\newcommand{\NlpTabRows}{}%
\DTLforeach*{nlpparams}{\prm=param,\vl=colvalue}{%
  \eappto\NlpTabRows{\expandonce{\prm} & \expandonce{\vl}\\\noexpand\hline}%
}%

\subsection{Языковые данные (иллюстративная конфигурация)}

Обучаемая надстройка для языковой серии задаётся конфигурациями репозитория: отдельно рекурсивный режим над замороженным Falcon и отдельно файнтюн через LoRA на том же чекпоинте (состав модулей и отсутствие совместного режима --- разд.~\ref{sec:slm_reuse}).

Основная серия экспериментов на естественном языке в полном объеме в данной работе не воспроизводится: ниже приведена \emph{типовая} постановка на публично доступных корпусах, чтобы зафиксировать масштабы и используемые показатели качества. Численные значения в табл.~\ref{tab:main_params} являются \textbf{ориентировочными заготовками} для масштаба величин и не должны трактоваться как результат конкретного численного эксперимента из кода проекта.

\begin{table}[H]
\centering
\caption{Иллюстративные параметры языкового эксперимента (публичные корпуса; числа условные; параметры задаются автором только как масштаб постановки; корпус WikiText-103 --- \cite{merity2016})}
\label{tab:main_params}
\begin{tabular}{|l|p{0.72\linewidth}|}
\hline
Параметр & Значение (иллюстрация) \\
\hline
\NlpTabRows
\end{tabular}
\end{table}

Для такой постановки на контрольной выборке ожидаются, например, перплексия порядка \(28\)--\(32\) и доля точного совпадения верхнего токена порядка \(0{,}35\)--\(0{,}42\) при фиксированной глубине рекурсии; эти диапазоны приведены только как ориентир методологии оценки и подлежат замене фактическими значениями при подключении языкового эксперимента к тому же коду, что использован для ARC-AGI.

\subsection{Сравниваемые варианты}

Сравнение должно быть устроено так, чтобы не смешивать вклад разных частей. Минимальный набор вариантов приведен в табл.~\ref{tab:model_variants}. Первая строка задает обычную модель с фиксированной глубиной. Остальные строки проверяют, что именно дает выигрыш: предобученная основа, рекурсивное уточнение, модель распределения времени или правило остановки.

\begin{table}[H]
\centering
\caption{Сравниваемые варианты модели}
\label{tab:model_variants}
\begin{tabularx}{\textwidth}{|l|Y|Y|Y|}
\hline
Вариант & Основа & Рекурсия & Время остановки \\
\hline
Фиксированная глубина & малая языковая модель & нет & не используется \\
\hline
Полный запуск рекурсии & малая языковая модель & до \(K\) шагов & всегда \(K\) \\
\hline
Ранняя остановка по энтропии & малая языковая модель & до \(K\) шагов & простая эвристика \\
\hline
Распределение \(\tau^\star\) & малая языковая модель & до \(K\) шагов & одно из пяти правил \\
\hline
LoRA~\cite{hu2021} & та же замороженная основа & нет (только малоранговые добавки) & не используется \\
\hline
\end{tabularx}
\end{table}

\subsection{Исследование вклада компонентов}

Исследование вклада компонентов нужно ограничить только тем, что относится к основной гипотезе. Проверяются четыре фактора: форма распределения времени, тип разметки, правило остановки и способ переиспользования малой языковой модели. Шаблон таблицы приведен в табл.~\ref{tab:ablation_template}.

\begin{table}[H]
\centering
\caption{Шаблон таблицы исследования вклада компонентов}
\label{tab:ablation_template}
\begin{tabularx}{\textwidth}{|l|Y|Y|Y|Y|}
\hline
Вариант & Потеря речи & \(\mathrm{NLL}_{\tau}\) & Среднее число шагов & Экономия \\
\hline
Конечное распределение + накопленная вероятность & \placeholder & \placeholder & \placeholder & \placeholder \\
\hline
Сглаженная разметка + накопленная вероятность & \placeholder & \placeholder & \placeholder & \placeholder \\
\hline
Смесь геометрических + малое будущее улучшение & \placeholder & \placeholder & \placeholder & \placeholder \\
\hline
Отрицательное биномиальное + квантиль & \placeholder & \placeholder & \placeholder & \placeholder \\
\hline
Бюджетное правило при заданном бюджете & \placeholder & \placeholder & \placeholder & \placeholder \\
\hline
\end{tabularx}
\end{table}

\subsection{Реализованный пайплайн ARC-AGI}\label{sec:arc_pipeline}

Для ARC-AGI зафиксирована серия из шести семейств распределения времени остановки. Для каждого семейства обучается рекурсивная модель с головой оракула; затем на контрольной выборке воспроизводится полная решётка правил остановки. В репозитории перечисление стратегий и порогов задано конфигурацией набора экспериментов (в коде~\texttt{configs/oracle\_sweep\_arc\_agi/}); сами стратегии имеют технические имена \texttt{cumulative\_probability}, \texttt{future\_improvement}, \texttt{hazard}, \texttt{quantile}, \texttt{budget}, а связка «режим формирования ответа / порог или бюджет» для каждого семейства подбирается по контрольной выборке. Режимы агрегирования кандидатов ответа фиксируются идентификаторами \texttt{last} и \texttt{argmax\_interval}. По итогам серии сохраняется сводная таблица; внутри каждого семейства строка выбирается по вспомогательному критерию \texttt{token\_acc\_best\_of\_policies} с режимом максимизации качества на контрольных данных.

Дополнительно для каждого базового обучения выполняется дополнительное обучение \emph{только} головы оракула (соответствующий запуск описан тем же кодовым деревом экспериментов); после этого перебор по решётке правил повторяется на тех же данных контроля, что исключает смещение сравнения из-за другой контрольной выборки.

Сводку по языковому моделированию на WikiText-103 см. разд.~\ref{sec:wikitext_results}. Серия на ARC-AGI и перебор по решётке стратегий с этапом донастройки головы оракула --- разд.~\ref{sec:arc_results}. Эмпирические распределения оракульного \(\tau^\star\) для ветки \texttt{finite\_discrete} на контроле приводятся в обоих результатных разделах.

